\section{注记}\label{sec:notes-differentiable-manifolds}

\subsection[第 \ref{sec:topological-manifolds} 节]{第~\ref{sec:topological-manifolds}~节}\label{subsec:notes-topological-manifolds}

先简要回顾所需的一般拓扑学的主要概念和结果 (详见 \RGcite{Mun00}).

\begin{enumerate}[label=(\arabic*),leftmargin=*]
\item 集合 $M$ 上的一个\textbf{拓扑}是满足下列性质的 $M$ 的子集族 $\mathcal{T}$:
\begin{enumerate}[label=(\roman*),leftmargin=*]
\item $\varnothing$ 和 $M$ 属于 $\mathcal{T}$;
\item $\mathcal{T}$ 的任意子族的成员之并属于 $\mathcal{T}$;
\item $\mathcal{T}$ 的任意有限子族的成员之交属于 $\mathcal{T}$.
\end{enumerate}
赋有拓扑 $\mathcal{T}$ 的集合 $M$ 称为\textbf{拓扑空间}. 若子集 $U\subset M$ 属于拓扑 $\mathcal{T}$, 则称它为 $M$ 的\textbf{开集}. 点 $p\in M$ 的一个\textbf{邻域}就是包含 $p$ 的开集 $U\in\mathcal{T}$. 一个\textbf{闭集} $F\subset M$ 是其补集 $M\setminus F$ 为开集的集合. 子集 $A\subset M$ 的\textbf{内部} $\operatorname{int}A$ 是包含在 $A$ 中的最大开集, 而其\textbf{闭包} $\overline{A}$ 是包含 $A$ 的最小闭集. 最后, $A\subset M$ 上的\textbf{子空间拓扑}是 $\mathcal{T}_A:=\{U\cap A\}_{U\in\mathcal{T}}$.
\item 如果对 $M$ 中每对不同的点 $p_1,p_2$ 都存在它们各自的邻域 $U_1,U_2$, 使得 $U_1\cap U_2=\varnothing$, 则称拓扑空间 $(M,\mathcal{T})$ 是 \textbf{Hausdorff} 的.
\item $M$ 上的拓扑 $\mathcal{T}$ 的一个\textbf{基}是一个子族 $\mathcal{B}\subset\mathcal{T}$, 使得对每个点 $p\in M$ 以及每个包含 $p$ 的开集 $U$, 都有一个基元素 $B\in\mathcal{B}$ 满足 $p\in B\subset U$. 若 $\mathcal{B}$ 是拓扑 $\mathcal{T}$ 的基, 则 $\mathcal{T}$ 的任意成员都是 $\mathcal{B}$ 的某些成员之并. 若 $\mathcal{T}$ 有可数基, 则称拓扑空间 $(M,\mathcal{T})$ 满足\textbf{第二可数公理}.
\item 两个拓扑空间之间的映射 $f:M\to N$ 称为\textbf{连续的}, 若对每个开集 $U\subset N$, 原像 $f^{-1}(U)$ 是 $M$ 的开子集. 若双射 $f$ 和它的逆 $f^{-1}$ 都连续, 则称其为\textbf{同胚}.
\item 拓扑空间 $(M,\mathcal{T})$ 的一个\textbf{开覆盖}是一个族 $\{U_\alpha\}\subset\mathcal{T}$, 满足 $\bigcup_\alpha U_\alpha=M$. \textbf{子覆盖}是仍为开覆盖的子族 $\{V_\beta\}\subset\{U_\alpha\}$. 若每个开覆盖都有有限子覆盖, 则称拓扑空间是\textbf{紧致的}. 若 $A\subset M$ 赋予子空间拓扑后是紧致拓扑空间, 则称 $A$ 为\textbf{紧子集}. 容易看出, 连续映射把紧集映为紧集.
\item 若 $M$ 中既开又闭的子集只有 $\varnothing$ 和 $M$, 则称拓扑空间是\textbf{连通的}. 若 $A\subset M$ 赋予子空间拓扑后是连通拓扑空间, 则称 $A$ 为\textbf{连通子集}. 容易看出, 连续映射把连通集映为连通集.
\item 设 $(M,\mathcal{T})$ 是拓扑空间. 若对 $p$ 的每个邻域 $V$ 都存在 $N\in\mathbb{N}$, 使得当 $n>N$ 时 $p_n\in V$, 则称 $M$ 中序列 $\{p_n\}$ \textbf{收敛}到 $p\in M$. 若 $(M,\mathcal{T})$ 是 Hausdorff 的, 收敛序列的极限唯一. 若 $(M,\mathcal{T})$ 还满足第二可数公理, 则 $F\subset M$ 是闭集, 当且仅当 $F$ 中每个收敛序列的极限都在 $F$ 中; $K\subset M$ 是紧集, 当且仅当 $K$ 中每个序列都在 $K$ 中有子列极限.
\item 若 $M$ 与 $N$ 是拓扑空间, 则 $M$ 的开子集与 $N$ 的开子集的所有 Cartesian 乘积, \phantomsection\label{idx:basis-for-product-topology}构成 $M\times N$ 上一个拓扑的基. 这一拓扑称为\textbf{乘积拓扑}. 注意, 在此拓扑下, 典范投影都是连续映射.
\item 集合 $M$ 上的一个\textbf{等价关系} $\sim$ 是具有下列性质的关系:
\begin{enumerate}[label=(\roman*),leftmargin=*]
\item \emph{自反性}: 对每个 $p\in M$, 有 $p\sim p$;
\item \emph{对称性}: 若 $p\sim q$, 则 $q\sim p$;
\item \emph{传递性}: 若 $p\sim q$ 且 $q\sim r$, 则 $p\sim r$.
\end{enumerate}
给定点 $p\in M$, 定义 $p$ 的\textbf{等价类}为集合
\[
  [p]=\{q\in M\mid q\sim p\}.
\]
由自反性可知 $p\in[p]$. 对集合 $M$ 上的任意等价关系 $\sim$, 相应的等价类所成的集合称为\textbf{商空间}, 记为 $M/{\sim}$. 存在一个典范投影 $\pi:M\to M/{\sim}$, 将 $M$ 的每个元素映到它的等价类. 若 $M$ 是拓扑空间, 可以在商空间上定义一个拓扑 (称为\textbf{商拓扑}): 子集 $V\subset M/{\sim}$ 是开集, 当且仅当 $\pi^{-1}(V)$ 在 $M$ 中是开集. 在这一拓扑下, $\pi$ 连续. 我们想知道某些商空间是否为 Hausdorff 空间. 为此, 下面的定义很有用.

\begin{RiemannDefinition}\label{def:open-equivalence-relation}
若映射 $\pi:M\to M/{\sim}$ 是开映射, 即对每个开集 $U\subset M$, 集合 $[U]:=\pi(U)$ 都是开集, 就称拓扑空间 $M$ 上的等价关系 $\sim$ 是\textbf{开的}.
\end{RiemannDefinition}

于是有

\begin{RiemannProposition}\label{prop:topology-of-quotients}
设 $\sim$ 是 $M$ 上的开等价关系, 令 $R=\{(p,q)\in M\times M\mid p\sim q\}$. 则商空间是 Hausdorff 的, 当且仅当 $R$ 在 $M\times M$ 中是闭的.
\end{RiemannProposition}

\begin{RiemannProof}
假设 $R$ 是闭的. 取 $[p],[q]\in M/{\sim}$ 且 $[p]\ne[q]$. 则 $p\not\sim q$, 且 $(p,q)\notin R$. 因为 $R$ 是闭的, 分别存在包含 $p,q$ 的开集 $U,V$, 使得 $(U\times V)\cap R=\varnothing$. 这蕴含 $[U]\cap[V]=\varnothing$. 事实上, 若存在一点 $[r]\in[U]\cap[V]$, 则 $r$ 分别与某些点 $p'\in U$ 和 $q'\in V$ 等价 (即 $p'\sim r$ 且 $r\sim q'$). 于是有 $p'\sim q'$ (从而 $(p',q')\in R$), 所以 $(U\times V)\cap R$ 非空. 因为 $[U]$ 和 $[V]$ 是开集 ($\sim$ 是开等价关系), 可知 $M/{\sim}$ 是 Hausdorff 的.

反过来, 假设 $M/{\sim}$ 是 Hausdorff 的. 若 $(p,q)\notin R$, 则 $p\not\sim q$ 且 $[p]\ne[q]$, 因而存在分别包含 $[p]$ 和 $[q]$ 的开集 $\widetilde U,\widetilde V\subset M/{\sim}$, 使 $\widetilde U\cap\widetilde V=\varnothing$. 集合 $U:=\pi^{-1}(\widetilde U)$ 与 $V:=\pi^{-1}(\widetilde V)$ 在 $M$ 中是开集, 并且 $(U\times V)\cap R=\varnothing$. 事实上, 若非如此, 就存在 $p'\in U$ 和 $q'\in V$ 使得 $p'\sim q'$. 于是有 $[p']=[q']$, 与 $\widetilde U\cap\widetilde V=\varnothing$ 矛盾 (因为 $[p']\in\pi(U)=\widetilde U$, 且 $[q']\in\pi(V)=\widetilde V$). 由于 $(p,q)\in U\times V\subset(M\times M)\setminus R$ 且 $U\times V$ 是开集, 可知 $(M\times M)\setminus R$ 是开集, 因而 $R$ 是闭的.
\end{RiemannProof}
\end{enumerate}

\subsection[第 \ref{sec:differentiable-manifolds} 节]{第~\ref{sec:differentiable-manifolds}~节}\label{subsec:notes-differentiability}

\begin{enumerate}[label=(\arabic*),leftmargin=*]
\item 先回顾 $\mathbb{R}^n$ 上映射的可微性的一些事实. 定义在 $\mathbb{R}^n$ 的开子集 $U$ 上的函数 $f:U\to\mathbb{R}$ 称为在 $U$ 上\textbf{连续可微}, 若所有偏导数 $\frac{\partial f}{\partial x^1},\ldots,\frac{\partial f}{\partial x^n}$ 在 $U$ 上存在且连续. 本书中, \emph{可微}和\emph{光滑}指\textbf{无穷次可微}, 即所有偏导数 $\frac{\partial^k f}{\partial x^{i^1}\cdots\partial x^{i^k}}$ 在 $U$ 上都存在且连续. 类似地, 定义在 $\mathbb{R}^n$ 的开子集上的映射 $f:U\to\mathbb{R}^m$ 称为\textbf{可微}或\textbf{光滑}, 若它的所有坐标函数 $f^i$ 也具有相同性质, 即都具有各阶连续偏导数. 若 $f$ 在 $U$ 上可微, 那么它在 $U$ 的每个点的\textbf{导数}是线性映射 $Df:\mathbb{R}^n\to\mathbb{R}^m$, 它在 $\mathbb{R}^n$ 和 $\mathbb{R}^m$ 的典范基下由 \textbf{Jacobian 矩阵}
\[
  Df=\begin{bmatrix}
    \dfrac{\partial f^1}{\partial x^1}&\cdots&\dfrac{\partial f^1}{\partial x^n}\\
    \vdots&&\vdots\\
    \dfrac{\partial f^m}{\partial x^1}&\cdots&\dfrac{\partial f^m}{\partial x^n}
  \end{bmatrix}.
\]
表示.

若存在开集 $U\supset A$ 和可微映射 $\widetilde f:U\to\mathbb{R}^m$ 使得 $f=\widetilde f|_A$, 则称定义在任意集合 $A\subset\mathbb{R}^n$ (不必是开集) 上的映射 $f:A\to\mathbb{R}^m$ \textbf{在 $A$ 上可微}.
\end{enumerate}

\subsection[第 \ref{sec:tangent-space} 节]{第~\ref{sec:tangent-space}~节}\label{subsec:notes-tangent-space}

\begin{enumerate}[label=(\arabic*),leftmargin=*]
\item 设 $E$, $B$ 和 $F$ 是光滑流形, $\pi:E\to B$ 是可微映射. 若满足以下条件, 则称 $\pi:E\to B$ 为一个以 $B$ 为\textbf{底空间}, 以 $E$ 为\textbf{全空间}, 以 $F$ 为\textbf{纤维}的\textbf{纤维丛}:
\begin{enumerate}[label=(\roman*),leftmargin=*]
\item 映射 $\pi$ 为满射;
\item 存在 $B$ 的开覆盖 $\{U_\alpha\}$ 以及微分同胚 $\psi_\alpha:\pi^{-1}(U_\alpha)\to U_\alpha\times F$, 使得对每个 $b\in U_\alpha$, 有 $\psi_\alpha(\pi^{-1}(b))=\{b\}\times F$.
\end{enumerate}
\end{enumerate}

\subsection[第 \ref{sec:immersions-embeddings} 节]{第~\ref{sec:immersions-embeddings}~节}\label{subsec:notes-inverse-function-theorem}

\begin{enumerate}[label=(\arabic*),leftmargin=*]
\item (\emph{反函数定理}) 设 $f:U\subset\mathbb{R}^n\to\mathbb{R}^n$ 是光滑函数, 且 $p\in U$ 使得 $(\mathrm{d}f)_p$ 是线性同构. 则存在包含 $p$ 的开子集 $V\subset U$, 使得 $f|_V:V\to f(V)$ 是微分同胚. 此外, 对所有 $q\in V$,
\[
  (\mathrm{d}(f|_V)^{-1})_{f(q)}=((\mathrm{d}(f|_V))_q)^{-1}.
\]
\end{enumerate}

\subsection[第 \ref{sec:lie-groups} 节]{第~\ref{sec:lie-groups}~节}\label{subsec:notes-proper-maps-and-fundamental-group}

\begin{enumerate}[label=(\arabic*),leftmargin=*]
\item \textbf{群}是赋有二元运算 $\cdot:G\times G\to G$ 的集合 $G$, 该运算满足:
\begin{enumerate}[label=(\roman*),leftmargin=*]
\item \textbf{结合律}: 对所有 $g_1,g_2,g_3\in G$, 有 $g_1\cdot(g_2\cdot g_3)=(g_1\cdot g_2)\cdot g_3$;
\item \textbf{单位元存在}: 存在元素 $e\in G$, 使得对所有 $g\in G$, 有 $e\cdot g=g\cdot e=g$;
\item \textbf{逆元存在}: 对所有 $g\in G$, 都有 $g^{-1}\in G$, 使得 $g\cdot g^{-1}=g^{-1}\cdot g=e$.
\end{enumerate}
若群运算可交换, 即对所有 $g_1,g_2\in G$ 有 $g_1\cdot g_2=g_2\cdot g_1$, 则称该群是 \textbf{Abel 群}. 若 $\cdot$ 到 $H\times H$ 的限制是 $H$ 上的二元运算, 且 $H$ 在这个运算下是群, 就称子集 $H\subset G$ 为 $G$ 的\textbf{子群}. 若对所有 $g\in G$, $h\in H$ 有 $ghg^{-1}\in H$, 则称子群 $H\subset G$ 是\textbf{正规的}. 若两个群 $G$ 和 $H$ 之间的映射 $f:G\to H$ 对所有 $g_1,g_2\in G$ 满足 $f(g_1\cdot g_2)=f(g_1)\cdot f(g_2)$, 则称它为\textbf{群同态}. \textbf{同构}是双射的同态. 群同态 $f:G\to H$ 的\textbf{核}是子集 $\ker(f)=\{g\in G\mid f(g)=e\}$, 容易看出它是 $G$ 的正规子群.
\item 设 $M$ 与 $N$ 是拓扑流形. 若每个紧集 $K\subset N$ 的原像 $f^{-1}(K)$ 是紧的, 就称映射 $f:M\to N$ 为\textbf{固有的}. 若 $f$ 还连续, 则 $f$ 是\textbf{闭映射}, 即它将闭集映为闭集. 为说明这一点, 取闭集 $F\subset M$, 考虑 $f(F)$ 中的收敛序列 $\{q_n\}$, 其中 $q_n\to q$. 容易看出, 集合 $\{q_n\mid n\in\mathbb{N}\}$ 的闭包 $K$ 是紧的. 由于 $f$ 固有, $f^{-1}(K)$ 也是紧的. 对每个 $n\in\mathbb{N}$, 选取 $p_n\in F$ 使得 $f(p_n)=q_n$. 则 $p_n\in f^{-1}(K)$, 因而 $\{p_n\}$ 必有一个子列极限 $p\in F$ (因为 $F$ 是闭的). 若 $\{p_{n_k}\}$ 是收敛到 $p$ 的一个子序列, 则 $q_{n_k}=f(p_{n_k})\to f(p)$ (因为 $f$ 连续). 因此 $q=f(p)\in f(F)$, 所以 $f(F)$ 是闭的.
\item 设 $f,g:X\to Y$ 是拓扑空间之间的两个连续映射, 令 $I=[0,1]$. 若存在连续映射 $H:I\times X\to Y$, 使得对每个 $x\in X$, 有 $H(0,x)=f(x)$ 且 $H(1,x)=g(x)$, 就称 $f$ 与 $g$ \textbf{同伦}. 映射 $H$ 称为一个\textbf{同伦}.

映射的同伦在 $X$ 到 $Y$ 的连续映射全体上构成等价关系. 作为应用, 在流形 $M$ 上固定一个\textbf{基点} $p$, 考虑满足 $f(0)=f(1)=p$ 的连续映射 $f:I\to M$ (这些映射称为\textbf{以 $p$ 为基点的环路}) 的同伦类, 其中另加条件: 对所有 $t\in I$, $H(t,0)=H(t,1)=p$. 这些同伦类所成的集合称为 $M$ \textbf{相对于基点 $p$ 的基本群}, 记为 $\pi_1(M,p)$. 它的元素中有\textbf{以 $p$ 为基点的常值环路}的类, 该环路由对每个 $t\in I$ 取 $f(t)=p$ 给出. 注意, 集合 $\pi_1(M,p)$ 在运算 $*$ (\textbf{环路的复合}) 下确实是群. 定义 $[f]*[g]:=[h]$, 其中 $h:I\to M$ 由
\[
  h(t)=\begin{cases}
    f(2t)&\text{若 }t\in[0,\frac12]\\
    g(2t-1)&\text{若 }t\in[\frac12,1]
  \end{cases}.
\]
给出. 这个群的单位元是以 $p$ 为基点的常值环路的等价类.

若 $M$ 连通且它是 $\pi_1(M,p)$ 中唯一的类, 则称 $M$ 是\textbf{单连通的}. 这意味着通过 $p$ 的每条环路都能连续变形为常值环路. 这一性质与点 $p$ 的选取无关, 并且等价于 $M$ 中任何闭路径都能连续变形为常值环路.
\item \textbf{四元数}是 Hamilton 在 1843 年引入的复数的推广. 他当时考虑形如 $a+bi+cj+dk$ 的数, 其中 $a,b,c,d\in\mathbb{R}$, 且
\[
  i^2=j^2=k^2=ijk=-1.
\]
形式上, 四元数集合 $\mathbb{H}$ 就是 $\mathbb{R}^4$, 其中
\begin{align*}
  1&=(1,0,0,0)\\ i&=(0,1,0,0)\\
  j&=(0,0,1,0)\\ k&=(0,0,0,1)
\end{align*}
并赋有由 Hamilton 公式 (以及 $1$ 为单位元的假设) 确定的双线性结合乘法. 在这些定义下, $\mathbb{H}$ 是一个除环, 即 $(\mathbb{H}\setminus\{0\},\cdot)$ 是 (非交换) 群, 且乘法对加法满足分配律.

四元数 $a+bi+cj+ik$ 的\textbf{实部}是 $a$, 而它的\textbf{向量部分}是 $bi+cj+dk$. 向量部分为零的四元数与实数等同, 实部为零的四元数与 $\mathbb{R}^3$ 中的向量等同. 四元数的\textbf{范数}是通常的 Euclidean 范数.
\end{enumerate}

\subsection{文献说明}\label{subsec:bibliographical-notes-differentiable-manifolds}

本章的材料完全是标准内容, 几乎可以在任何一本微分几何书中找到 (例如 [\RGref{bib:Boo03}, \RGref{bib:dC93}, \RGref{bib:GHL04}]). 浸入和嵌入是\textbf{微分拓扑}的起点, 可参见 [\RGref{bib:GP73}, \RGref{bib:Mil97}]. Lie 群和 Lie 代数是数学的一个广阔领域, 本书无法充分展开. 例如可参见 [\RGref{bib:BtD03}, \RGref{bib:DK99}, \RGref{bib:War83}]. 关于基本群和覆盖空间的更多细节, 例如可参见 \RGcite{Mun00}.
